
%---------------------------------------------------------------------------
% Custom Definitions for User Name, Course Code and Title, Instructor, etc.
%
\newcommand{\thiscoursecode}{AM 207}
\newcommand{\thiscoursename}{Advanced Scientific Computing}
\newcommand{\thiscoursenameb}{Stochastic Methods for Data Analysis, Inference and Optimization}
\newcommand{\thisprof}{ Verena Kaynig-Fittkau $\&$ Pavlos Protopapas}
\newcommand{\thistf}{Giri Gopalan}
\newcommand{\me}{Karen Yu, Nick Vasios, Thibaut Perol}
\newcommand{\thisterm}{Spring 2016}
%---------------------------------------------------------------------------
%###########################################################################
%---------------------------------------------------------------------------
%                           Code for the Front Page
%
\newcommand{\notefront} {
\begin{center}
    \includegraphics[width=2cm]{SEAS_S}\\[8pt]
    \textbf{\Huge{\thiscoursecode}}{\Huge \\[8pt]}
    {\large{\noun{\thiscoursename}}}\\[8pt]
    {\large{\noun{\thiscoursenameb}}}\\[8pt]
    {{\noun {Harvard John A. Paulson School of Engineering and Applied Sciences}}}
\end{center}
}
%---------------------------------------------------------------------------
%###########################################################################
%---------------------------------------------------------------------------


\newcommand{\no}{\noindent}
\newcommand{\pr}{\partial}
%
\def\beq#1\eeq{\begin{equation}#1\end{equation}}
\def\beqa#1\eeqa{\begin{eqnarray}#1\end{eqnarray}}
\def\beqas#1\eeqas{\begin{eqnarray*}#1\end{eqnarray*}}

\def\blst#1\elst{\begin{lstlisting}#1\end{lstlisting}}



% Quick Underbrace (Math Mode)
\newcommand{\ub}[1]{\underbrace{#1}}

% Quick Underbrace with text (Text Mode)
\newcommand{\ubt}[2]{$\underbrace{\mbox{#1}}_{\mbox{#2}}$}

% Combination Notation (Choose n from p) (Math Mode)
\newcommand{\comb}[2]{{#1 \choose #2}}

% Easy Piecewise function
% 1: Case 1
% 2: Domain for case 1
% 3: Case 2
% 4: Domain for case 2
\newcommand{\piecewise}[4]{\left\{
     \begin{array}{lr}
       #1 & \mbox{#2} \\
       &\\ #3 & \mbox{#4}
     \end{array}
   \right.}

% Easy Slanted Text
\newcommand{\noun}[1]{\textsc{#1}}
% Easy Text TT
\newcommand{\tc}[1]{\texttt{#1}}

% Easy <math>
\newcommand{\sip}[2]{\langle #1 , #2 \rangle}


% Small Greek Letter Shortcuts
\renewcommand\a{\alpha}
\renewcommand\v{\beta}
\newcommand\g{\gamma}
\renewcommand\d{\delta}
\newcommand\e{\epsilon}
\newcommand\ve{\varepsilon}
\def\k{\kappa}
\renewcommand\l{\lambda}
\newcommand\ldot{\dot{\lambda}}
\newcommand\m{\mu}
\newcommand\s{\sigma}
\renewcommand\t{\tau}
\renewcommand\th{\theta}
\newcommand\z{\zeta}
\newcommand\om{\omega}
\newcommand\ups{\upsilon}

% Some Capital Greek Letter Shortcuts
\newcommand\D{\Delta}
\newcommand\F{\Phi}
\renewcommand\L{\Lambda}
\newcommand\Sig{\Sigma}
\def\Om{\Omega}

% Bold Symbols for Vectors and Tensors
\newcommand\baa{\boldsymbol{\alpha}}
\newcommand\bbe{\boldsymbol{\beta}}
\newcommand\bfdel{\boldsymbol{\delta}}
\newcommand\bee{\boldsymbol{\epsilon}}
\newcommand\bve{\boldsymbol{\varepsilon}}
\newcommand\bfeta{\boldsymbol{\eta}}
\newcommand\bkappa{\boldsymbol{\kappa}}
\newcommand\bmi{\boldsymbol{\mu}}
\newcommand\bfxi{\boldsymbol{\xi}}
\newcommand\bs{\boldsymbol{\sigma}}
\newcommand\bom{\boldsymbol{\omega}}
\newcommand\bedot{\dot{\boldsymbol{\epsilon}}}
\newcommand\bGamma{\bf\Gamma}
\newcommand\bDel{\bf\Delta}
\newcommand\bPi{\boldsymbol{\Pi}}
\newcommand\bSig{\boldsymbol\Sigma}
\newcommand\bPhi{\boldsymbol{\Phi}}
\newcommand\bOm{\boldsymbol{\Omega}}
\newcommand\bgrad{\boldsymbol{\nabla}}
\newcommand\cor{\hat{\bs}}
\newcommand\ba{{\bf a}}
\newcommand\bb{{\bf b}}
\newcommand\bc{{\bf c}}
\newcommand\bff{{\bf f}}
\newcommand\bd{{\bf d}}
\newcommand\be{{\bf e}}
\newcommand\bg{{\bf g}}
\newcommand\bh{{\bf h}}
\newcommand\bj{{\bf j}}
\newcommand\bk{{\bf k}}
\newcommand\bl{\boldsymbol{l}}
\newcommand\bm{{\bf m}}
\newcommand\bn{{\bf n}}
\newcommand\bp{{\bf p}}
\newcommand\br{{\bf r}}
\newcommand\bx{{\bf x}}
\newcommand\by{{\bf y}}
\def\bq{{\bf q}}
\newcommand\bss{{\bf s}}
\newcommand\bt{{\bf t}}
\newcommand\bu{{\bf u}}
\newcommand\bv{{\bf v}}
\newcommand\bw{\boldsymbol{w}}
\newcommand\bups{\boldsymbol{\upsilon}}
\newcommand\btau{\boldsymbol{\tau}}
\newcommand\bwv{\{\boldsymbol{w}^{*}\}}
\newcommand\bwt{\{\boldsymbol{w}\}^T}
\newcommand\bwvt{\{\boldsymbol{w}^{*}\}^T}
\newcommand\brho{\boldsymbol{\rho}}
\newcommand\blam{\boldsymbol{\lambda}}
\newcommand\bA{{\bf A}}
\newcommand\bB{{\bf B}}
\newcommand\bC{{\bf C}}
\newcommand\bD{{\bf D}}
\newcommand\bE{{\bf E}}
\newcommand\bF{{\bf F}}
\newcommand\bG{{\bf G}}
\newcommand\bH{{\bf H}}
\newcommand\bI{{\bf I}}
\newcommand\bJ{{\bf J}}
\newcommand\bK{{\bf K}}
\newcommand\bL{{\bf L}}
\newcommand\bM{{\bf M}}
\newcommand\bN{{\bf N}}
\newcommand\bP{{\bf P}}
\newcommand\bQ{{\bf Q}}
\newcommand\bR{{\bf R}}
\newcommand\bS{{\bf S}}
\newcommand\bT{{\bf T}}
\newcommand\bU{{\bf U}}
\newcommand\bV{{\bf V}}
\newcommand\bW{{\bf W}}
\newcommand\bX{{\bf X}}
\newcommand\bY{{\bf Y}}
\newcommand\bZ{{\bf Z}}
\newcommand\bOM{{\bf \Omega}}
\newcommand\bPsi{{\bf \Psi}}
\newcommand\bLam{{\bf \Lambda}}
\newcommand\bzero{{\bf 0}}
% Other
\newcommand\Ccal{\mathcal{C}}
\newcommand\Dcal{\mathcal{D}}
\newcommand\Ical{\mathcal{I}}
\newcommand\bLcal{\boldsymbol{\mathcal{L}}}
\newcommand\bMcal{\boldsymbol{\mathcal{M}}}
\newcommand\bKcal{\boldsymbol{\mathcal{K}}}
\newcommand\bJcal{\boldsymbol{\mathcal{J}}}
\newcommand\bCcal{\boldsymbol{\mathcal{C}}}
\newcommand\bIcal{\boldsymbol{\mathcal{I}}}
\newcommand\bDcal{\boldsymbol{\mathcal{D}}}
\newcommand\Lcal{\mathcal{L}}
\newcommand\Mcal{\mathcal{M}}


\newcommand{\union}{\cup}
\newcommand{\inter}{\cap}

% Very useful shortcuts for sets
\newcommand{\nset}[1]{#1_1,\ldots,#1_n}
\newcommand{\setk}[2]{#1_1,\ldots,#1_#2}
\newcommand{\bnset}[1]{\{#1_1,\ldots,#1_n\}}
\newcommand{\bset}[2]{\{#1_1,\ldots,#1_#2\}}
\newcommand{\ton}[1]{#1=1,\ldots,n}
\newcommand{\tok}[2]{#1=1,\ldots,#2}

% Operators
\newcommand{\disc}{\operatorname{disc}}
\newcommand{\spn}{\operatorname{Span}}
\newcommand{\rank}{\operatorname{rank}}
\newcommand{\proj}{\operatorname{proj}}
\newcommand{\prp}{\operatorname{perp}}
\newcommand{\tran}[1]{#1^\mathrm{T}}

% Underlines
\newcommand{\ux}{\underline{x}}
\newcommand{\ua}{\underline{a}}
\newcommand{\uu}{\underline{u}}

\theoremstyle{definition}
%\newtheorem{cor}{Corollary}[section]
\newtheorem*{note}{Note}
\newtheorem{notation}{Notation}[section]
\newtheorem{exercise}{Exercise}[section]
\newtheorem*{solution}{Solution}
\newtheorem*{ans}{Answer}

\newcommand{\Lp}{\mathcal{L}^p}

\declaretheorem[name=Theorem,
refname={theorem,theorems},
Refname={Theorem,Theorems},
numberwithin=section]{thrm}

\declaretheorem[name=Lemma,
refname={lemma,lemmas},
Refname={Lemma,Lemmas},
numberwithin=section]{lem}

\declaretheorem[name=\color{crimson} Example,
refname={example,examples},
Refname={Example,Example},
numberwithin=section]{exmp}

\declaretheorem[name=Remark,
refname={remark,remarks},
Refname={Remark,Remarks},
numberwithin=section]{rem}

\declaretheorem[name=Proposition,
refname={proposition,propositions},
Refname={Proposition,Propositions},
numberwithin=section]{prop}

\declaretheorem[name=Definition,
refname={definition,definitions},
Refname={Definition,Definitions},
numberwithin=section]{defn}

\newcommand{\Python}
{
\begin{flushright}
  \textit{Python Example}
\end{flushright}\vspace{-2ex}
}
\newcommand{\var}{\mathrm{var}}
\newcommand{\cov}{\mathrm{cov}}

\DeclareMathOperator*{\Assembly}{\scalerel*{\mathrm{A}}{\Big[}}
\newcommand*\tick{\item[\Checkmark]} 

\usepackage{etoolbox}
\makeatletter
\patchcmd{\chapter}{\if@openright\cleardoublepage\else\clearpage\fi}{}{}{}
\makeatother